\documentclass{beamer}
%\documentclass[aspectratio=169]{beamer} % 使用 16:9 宽屏比例
\usepackage{hyperref}
\usepackage[T1]{fontenc}

% other packages
\usepackage{latexsym,amsmath,xcolor,multicol,booktabs,calligra}
\usepackage{graphicx,pstricks,listings,stackengine}

\author{Stella Rovazzi}
\title{Perturbation in XAI}
\subtitle{Understanding their affect on models and metrics}
\institute{University of Guelph}
\date{Summer 2026}
\usepackage{CNU}

% defs
\def\cmd#1{\texttt{\color{red}\footnotesize $\backslash$#1}}
\def\env#1{\texttt{\color{blue}\footnotesize #1}}
\definecolor{deepblue}{rgb}{0,0,0.5}
\definecolor{deepred}{rgb}{0.6,0,0}
\definecolor{deepgreen}{rgb}{0,0.5,0}
\definecolor{halfgray}{gray}{0.55}

\lstset{
    basicstyle=\ttfamily\small,
    keywordstyle=\bfseries\color{deepblue},
    emphstyle=\ttfamily\color{deepred},    % Custom highlighting style
    stringstyle=\color{deepgreen},
    numbers=left,
    numberstyle=\small\color{halfgray},
    rulesepcolor=\color{red!20!green!20!blue!20},
    frame=shadowbox,
}

\AtBeginSubsection{}
\begin{document}


\begin{frame}
    \titlepage
\end{frame}

{\AtBeginSection{}
\section{Introduction}}
\subsection{Research Objectives and Expected Outcomes}

\begin{frame}{Introduction} % Stella

    \textbf{Motivation:} Help researchers and developers make informed perturbation choices to improve explanations. \\
        \vspace{0.5cm}      
   
    \begin{enumerate}
        \item To examine how the choice of perturbation technique effects an explanation's output (LIME)
        \item To examine how the choice of perturbation technique effects a metric's output (Sensitivity, Fidelity, Relative Stability)
        \item To explore common patterns between the experiment's tests and understand the effects of the studied perturbation techniques
    \end{enumerate}
        \vspace{0.5cm}      
	\textbf{Expected Outcome:} Resampling, or recreation, is the most reliable and universal perturbation technique.

\end{frame}


{\setbeamertemplate{subsection page}{}
\subsection{Methodology}} % Stella

\begin{frame} % Stella or Adil: Lattice Definition and why secure --> Fixing noise from lengthy lattices --> Making algebraic w LWE --> Putting togetrher
\tableofcontents[sectionstyle=show,subsectionstyle=show/shaded/hide,subsubsectionstyle=show/shaded/hide]
\end{frame}

\begin{frame}{Methodology - Model}
            \begin{enumerate}
                \item Train one-hidden layer FNN on benchmark datset
                \item Fit the LIME model with each of the perturbation techniques
                \item Calculate the gradient (to find the "true" feature importance) of the black-box
                \item Evaluate the explanatory vectors' distance from the "true" explanatory vector with the normalized euclidean distance, cosine similarity, and spearman ranking
            \end{enumerate}
\end{frame}

\begin{frame}{Methodology - Metric}
        \begin{enumerate}
            \item Train one-hidden layer FNN model on benchmark datset
            \item Fit decision tree explainer(s) to each output of FNN model (control local models)
            \item Find the Fidelity of the model with each of the perturbation techniques
            \item Compare the different results of the metrics, dependent on the perturbation chosen
            \item Repeat steps 1 to 4 for sensitivity
        \end{enumerate}
\end{frame}

\begin{frame}{Perturbation Methods}
    \begin{itemize}
        \item \textbf{Gaussian} - Resampling method following a $N(0, 1)$, shifted to centre around the feature's mean
        \item \textbf{Latin Hypercube Sampling} - Another resampling method, with a different distribution
        \item \textbf{Replace with Mean}* - Replaces given value with the mean
        \item \textbf{Replace with Mean conditioned on threshold}* - Replace given value with mean if it follows some condition 
        \item \textbf{Reduction}* - Replace value with a minimized amount (based on softmasking)
    \end{itemize}

    \vspace{0.5cm}

    \textit{* are custom methods that also use a binary masks to ensure the entire feature column is not replaced by another value}
\end{frame}

\begin{frame}{Expiremental Setting}

\tiny
\setlength{\tabcolsep}{3pt}
\renewcommand{\arraystretch}{0.9}
\resizebox{\textwidth}{!}{      
    \begin{tabular}{|p{0.22\textwidth}|p{0.32\textwidth}|p{0.28\textwidth}|}
        \hline
        \textbf{Research question} &
        \textbf{Independent variable} &
        \textbf{Dependent variable} \\
        \hline
        Affect on Models &
        \begin{itemize}
            \item Dataset (Wine Data)
            \item Feed-Forward Neural Network
            \item "True" Explantory Vector (or calculated gradient)
        \end{itemize} &
        Lime's Explanatory Vectors \\
        \hline 
        Affect on Metrics &
        \begin{itemize}
            \item Dataset (Wine Data)
            \item FNN
            \item Decision Tree Explainers (One for each observation of test set)
            \item Perturbed inputs (metrics calculated in one loop and seed set)
        \end{itemize} &
        Metric Values (Sensitivty, Infidelity) \\
        \hline
    \end{tabular}
    }
\end{frame}

\section{Results}
{\setbeamertemplate{subsection page}{}
\subsection{LIME}} % Stella

\begin{frame}{Normalized Euclidean Distance}
   \begin{figure}% [H] forces the image to stay exactly HERE
    \includegraphics[width=0.85\textwidth]{euclidean_vs_true} % Omit extension (.jpg/.png)
    \caption{Comparing the gradient to LIME perturbed explanatory vectors.}
    \end{figure}
\end{frame}

\begin{frame}{Cosine Similarity}
   \begin{figure}% [H] forces the image to stay exactly HERE
    \includegraphics[width=0.85\textwidth]{cosine_similarity_vs_true} % Omit extension (.jpg/.png)
    \caption{Comparing the gradient to LIME perturbed explanatory vectors.}
    \end{figure}
\end{frame}

\begin{frame}{Spearman Ranking}
   \begin{figure}% [H] forces the image to stay exactly HERE
    \includegraphics[width=0.85\textwidth]{spearman_rank_vs_true} % Omit extension (.jpg/.png)
    \caption{Comparing the gradient to LIME perturbed explanatory vectors.}
    \end{figure}
\end{frame}

{\setbeamertemplate{subsection page}{}
\subsection{Infidelity}} % Stella

\begin{frame}{Infidelity}
    \begin{figure}% [H] forces the image to stay exactly HERE
    \includegraphics[width=0.85\textwidth]{tree_infid_class_both} % Omit extension (.jpg/.png)
    \caption{Infidelity of Each Method}
    \end{figure}
\end{frame}


\begin{frame}{Infidelity}
    \begin{figure}% [H] forces the image to stay exactly HERE
    \includegraphics[width=0.85\textwidth]{ActvsPred.png} % Omit extension (.jpg/.png)
    \caption{Comparison of Assessing the Change}
    \end{figure}
\end{frame}

{\setbeamertemplate{subsection page}{}
\subsection{Sensitivity}} % Stella

\begin{frame}{Sensitivity}
    \begin{figure}% [H] forces the image to stay exactly HERE
    \includegraphics[width=0.85\textwidth]{maximum_sensitivity} % Omit extension (.jpg/.png)
    \end{figure}
\end{frame}


\begin{frame}{Sensitivity}
    \begin{figure}% [H] forces the image to stay exactly HERE
    \includegraphics[width=0.85\textwidth]{mean_sensitivity} % Omit extension (.jpg/.png)
    \end{figure}
\end{frame}


\begin{frame}{Sensitivity}
    \begin{figure}% [H] forces the image to stay exactly HERE
    \includegraphics[width=0.85\textwidth]{worstvsmax} % Omit extension (.jpg/.png)
    \end{figure}
\end{frame}

\section{Findings}

\begin{frame}{Findings/Current Thoughts}
    \begin{itemize}
        \item Some form of resampling is likely the strongest for an explainer model, over other perturbations. You want to mimic the data's behaviour
        properly.
        \item Reduction can force a perturbation based metric make a model look better than it may be
        \item I have the decision tree's comparisons to the gradient vector - and I feel like that may make my metric section stronger but
        I haven't figured out how to properly tie it in yet.
    \end{itemize}
\end{frame}
\end{document}
